

\section{Case Studies}

%In this section, we summarize the published studies that we have developed using the DPL concept from two perspectives: scalability by combining multiple motions, and application to complex manipulation tasks.
This section summarizes the published studies we developed using the DPL concept from two perspectives: scalability by combining multiple motions and application to complex manipulation tasks.



\subsection{Compound Generation of Robot Motions}
\label{sec_case_studies}

%As mentioned in the first section, approaches based on large-scale data and trials have high training costs, so it is necessary to perform a wide range of tasks with limited robot motion data. Here, since many practical tasks are essentially long-horizon tasks that consist of multiple steps, an approach that expresses them by dividing and combining short tasks (subtasks) is effective~\cite{Sacerdoti1974ai,Wolpert1998,Schaal2003isrr}. This kind of "compound generation of robot motions" has the following advantages.
As mentioned in the first section, approaches based on large-scale data and trials have high training costs; therefore, it is necessary to perform a wide range of tasks using limited robot motion data. It is an effective approach to express practical tasks by dividing and combining short tasks (subtasks)because many are essentially long-horizon tasks consisting of multiple steps~\cite{Sacerdoti1974ai,Wolpert1998,Schaal2003isrr}. This type of "compound generation of robot motions" has the following advantages:
\begin{itemize}
    \item Using subtasks included in a long-term task for another task (module reuse)
    \item Executing new tasks by combining subtasks (extending functionality)
\end{itemize}
%With the above background, this section discusses the elements required to combine multiple motions and introduces robot applications with the DPL concept.
With this background, this section discusses the elements required to combine multiple motions and introduces robot applications using the DPL concept.


\subsubsection{Multi-Module Switching based on Prediction Error}

%One of the simple approaches to compound-generating robot motions is the cooperation of modularized functions~\cite{Haruno2001,Wolpert1998}. Generally, a robot system consists of a group of multiple modules, and the output is determined by combining them. Even in robot learning, it can be assumed that by combining or switching models with different functions in parallel, the overall task performance can be expanded without increasing the computational load. Furthermore, by having a distributed system configuration, it is possible to continuously add functionality by adding modules.
A simple approach to compound-generating robot motion involves the cooperation of modularized functions~\cite{Haruno2001,Wolpert1998}. Generally, a robotic system consists of multiple modules, and the output is determined by combining them. Even in robot learning, it can be assumed that by combining or switching models with different functions in parallel, the overall task performance can be improved without increasing the computational load. Furthermore, with a distributed system configuration, it is possible to continuously add functionality by adding modules.


%One method to achieve the above in robot learning is to switch subtask modules based on the prediction error of the model. As mentioned in the previous Section, RNNs can retain various dynamics obtained from the learning process as unique internal states. When generating robot motion, RNN predicts from the internal state and current sensorimotor information, and, prediction errors tend to increase in untrained states. Therefore, when the prediction error increases during motion generation, the RNN has the property of switching from the dynamics to minimize the prediction error, and this can be used to switch modules. When a large prediction error that cannot be handled in real-time occurs, it cannot be handled by the state transition of the attracter embedded in individual RNNs. In that case, it is necessary to adjust the overall sensory prediction and motion generation by coordinating multiple modules according to the amount of the prediction error.
One method for achieving this in robot learning is to switch the subtask modules based on the model's prediction error. As mentioned in the previous Section, RNNs can retain the various dynamics obtained from the learning process as unique internal states. When generating robot motion, RNN predicts from the internal state and current sensorimotor information, and the prediction errors tend to increase in untrained states. Therefore, when the prediction error increases during motion generation, the RNN switches from dynamics to minimize the prediction error, which can be used to switch modules. When a large prediction error that cannot be handled in real-time occurs, it cannot be handled by the state transition of the attracter embedded in the individual RNNs. In this case, it is necessary to adjust the overall sensory prediction and motion generation by coordinating multiple modules according to the prediction error.


%\cite{Ito2022sr} proposed a method that takes actions to reduce prediction errors in door-passing tasks by introducing multiple DPL models and a new mechanism that switches them according to prediction errors (Figure~\ref{fig_doorrobot}). The model was trained for each subtask, such as ``opening a door'' or ``passing through,'' and then combined these models to perform a series of tasks. Each model predicted the near future camera image from sensor information. By comparing it with the real camera image, the method calculated the change in confidence value, which indicates how accurately the robot can work in the real environment, over time. Furthermore, the robot executed actions appropriate to the situation by autonomously selecting the model with the highest confidence. Since the robot performed these calculations in real-time, it could handle complex tasks, without having to accurately design the timing of switching motions or the flow of motion. As a result of real robot experiments, we confirmed that it is possible to generate tasks in response to unlearned door positions and changes in door panel color and pattern\footnote{https://youtu.be/pdOA-PeYO9Q}.
\cite{Ito2022sr} proposed a method that takes actions to reduce prediction errors in door-passing tasks by introducing multiple DPL models and a new mechanism that switches them according to prediction errors (Figure~\ref{fig_doorrobot}). The model was trained for each subtask, such as ``opening a door'' or ``passing through,'' and then combined these models to perform a series of tasks. Each model predicts a near-future camera image from the sensor information. The method calculates the change in the confidence value, which indicates how accurately the robot can work in a real environment over time by comparing it with a real camera image. Furthermore, the robot executes actions that are appropriate for the situation by autonomously selecting the model with the highest confidence. It could handle complex tasks without having to accurately design the timing of the switching motions or the flow of motion because the robot performed these calculations in real time. Based on real robot experiments, we confirmed that generating tasks in response to unlearned door positions and changes in the door panel color and pattern\footnote{https://youtu.be/pdOA-PeYO9Q} is possible.


%-------------------- Fig --------------------
\begin{figure}[t]
\centering
\includegraphics[width=73mm]{DoorRobotv2.eps}
\caption{
Real-time switching of multiple motion generation models based on prediction errors~\cite{Ito2022sr}.
}
\label{fig_doorrobot}
\hspace{10mm}
\end{figure}


%On the other hand, \cite{Suzuki2021ral} worked on responding to robot states that are not included in the training data by utilizing multiple control modules with different characteristics. This is an attempt to deal with the extrapolation of motion datasets (undefined behaviors) that cannot be predicted by a learning-based controller. In the proposed method, we prepared a model-based controller that compensates the output of the RNN. And, we appropriately switched models according to the error between the predicted trajectory and the actual trajectory. Assuming that the RNN embedded the most recent information necessary for task execution in its internal state, we trained additional tasks to predict past trajectories from the internal state. By adding a new neuron that predicts additional tasks to the middle layer of the RNN, it is possible to add error recovery functionality without affecting the performance of the main task learned. In real robot experiments, we confirmed that by combining this method with the dynamic design of the RNN internal state described in the next section, it is possible to perform error recovery in response to multiple types of disturbances (environment contact and changes) from collaborative workers.
Contrarily, \cite{Suzuki2021ral} worked on responding to robot states not included in the training data by utilizing multiple control modules with different characteristics. This approach attempted to extrapolate motion datasets (undefined behaviors) that cannot be predicted using a learning-based controller. In the proposed method, we prepared a model-based controller that compensated for the output of the RNN. In addition, the models were appropriately switched based on the error between the predicted and actual trajectories. Assuming that the RNN embeds the most recent information necessary for task execution in its internal state, we train additional tasks to predict past trajectories from the internal state. By adding a new neuron that predicts additional tasks to the middle layer of the RNN, an error recovery functionality can be added without affecting the performance of the main learned task. In real robot experiments, we confirmed that, by combining this method with the dynamic design of the RNN internal state described in the next section, it is possible to recover errors in response to multiple types of disturbances (environmental contact and changes) from collaborative workers.


%In general robot systems, it is necessary to separately program responses to unexpected situations, which requires a lot of development cost and adjustment work. Additionally, when an unexpected disturbance occurs in the work environment, the robot has to take time to recognize the situation and replan, which causes the robot to move slowly. Our proposed methods~\cite{Ito2022sr,Suzuki2021ral} can respond to sudden situational changes and changes in operating procedures by switching multiple controllers according to prediction errors, and are expected to solve the above problems.
In general robot systems, it is necessary to separately program responses to unexpected situations, which require significant development costs and adjustment work. In addition, when an unexpected disturbance occurs in the work environment, the robot must take time to recognize the situation and replan, which causes it to move slowly. Our proposed methods~\cite{Ito2022sr,Suzuki2021ral} can respond to sudden situational changes and changes in operating procedures by switching multiple controllers according to prediction errors and are expected to solve the above problems.


\subsubsection{Switching of Multiple Behaviours by Designing the Dynamical Systems}

%Next, we discuss the dynamical system of the RNN as a method for handling local composite generation within a single module. The method described in the previous section focused on switching between modules, however, considering scalability, it is desirable to have a mechanism that can also switch operations within individual modules. This is because when a task is further subdivided, each module is required to handle multiple subtasks~\cite{Nair2019iclr,Yu2019iros}. Previous studies~\cite{Yin2021iclr,Garrett2021cras} have also investigated motion planning that combines subtasks. To connect multiple robot motions, it is necessary to express their dynamics in a combinable format within a module. By selecting appropriate dynamics of the RNN according to the sensory input, it becomes possible to execute long-term task operations in response to environmental changes. We implemented the above by explicitly embedding attractor dynamics in subtask units into the RNN.
Next, we discuss the dynamic system of the RNN as a method for handling local composite generation within a single module. The method described in the previous section focuses on switching between modules; however, considering scalability, it is desirable to have a mechanism that can switch operations within individual modules. This need arises because when a task is subdivided, each module must handle multiple subtasks~\cite{Nair2019iclr,Yu2019iros}. Previous studies~\cite{Yin2021iclr,Garrett2021cras} also investigated motion planning that combines subtasks. Expressing their dynamics in a combined format within a module is necessary to connect multiple robot motions. Long-term task operations can be executed in response to environmental changes by selecting the appropriate dynamics of the RNN according to the sensory input. We implemented the above by explicitly embedding the attractor dynamics in the subtask units into an RNN.


%In \cite{Kase2018icra}, to combine multiple subtasks, we intentionally embedded an attractor structure inside the model by using a loss function in the internal state of the RNN (Figure~\ref{fig_kase_attractor}). By adding a constraint loss to have a shared state that can be switched between the attractors of each subtask, the attractor switches its dynamics to reduce the prediction error. In the experiment, by combining multiple subtasks for a Put-in-Box task, we achieved motion generation in an arbitrary order of subtasks\footnote{https://youtu.be/QwogmDqOxfE}. Furthermore, we expanded the loss function of the internal state of RNN and designed an arbitrary attractor structure to increase the degree of freedom in switching postures~\cite{Kase2019iros}.
In \cite{Kase2018icra}, to combine multiple subtasks, we intentionally embedded an attractor structure inside the model using a loss function in the internal state of the RNN (Figure~\ref{fig_kase_attractor}). The attractor switches its dynamics to reduce the prediction error by adding a constraint loss to obtain a shared state that can be switched between the attractors of each subtask. In this experiment, combining multiple subtasks for a put-in-box task generated motions in an arbitrary order of subtasks\footnote{https://youtu.be/QwogmDqOxfE}. Furthermore, we expanded the loss function of the internal state of RNN and designed an arbitrary attractor structure to increase the degrees of freedom in switching postures~\cite{Kase2019iros}.


%-------------------- Fig --------------------
\setlength\textfloatsep{5pt}
\begin{figure}[t]
\centering
\includegraphics[width=\columnwidth]{icra2018_attractor_v2.eps}
\caption{
Embedding attractor dynamics into the internal state of RNN~\cite{Kase2018icra}.
}
\label{fig_kase_attractor}
\hspace{10mm}
\end{figure}


%\cite{Suzuki2018ral} introduced a hierarchical structure of task instruction input and RNN to learn how to switch subtasks during a long-horizon task. By inputting a task instruction vector at the designed point attractor, it is possible to respond to motion switching in ambiguous situations that cannot be determined based on sensor information alone. In the experiment, we verified a clothes-folding task that consisted of multiple subtasks and confirmed that the robot could perform according to the instructions. Furthermore, in the internal state of the RNN, instructions were embedded in the long-term time scale layer, and motion was embedded in the short-term time scale layer as an attractor structure. This indicated that the task context, language, and image input were appropriately structured to combine.
\cite{Suzuki2018ral} introduced a hierarchical structure of task instruction input and RNN to learn how to switch subtasks during a long-horizon task. Inputting a task instruction vector to the designed point attractor enables the response to motion switching in ambiguous situations that cannot be determined based on sensor information alone. In the experiment, we verified a clothes-folding task that consisted of multiple subtasks and confirmed that the robot could perform it according to the instructions. Furthermore, in the internal state of the RNN, instructions were embedded in the long-term timescale layer, and motion was embedded in the short-term timescale layer as an attractor structure. These embeddings indicated that the task context, language, and image inputs were appropriately combined.


%Although these motion combination/switching methods using attractor structures require the design of switching postures for subtasks, they can incorporate the human's intentions to some extent into the RNN output. Combined with the grounding with the language described below, it is expected that the robot will be able to switch its motions more flexibly.
Although these motion combination/switching methods using attractor structures require the design of switching postures for subtasks, they can incorporate human intentions to some extent into the RNN output. The robot is expected to be able to switch its motions more flexibly, combined with grounding, using the language described below.



\subsubsection{Language-Conditioned Motion Generation}

%To express more primitive combinations of actions, it is necessary to ground each segmented subtask and language expression. By learning that the given language instructions and individual tasks are semantically grounded, the RNN expresses the internal dynamics with the semantic and syntactic structure of the language system as a type of constraint. This internal representation can be extended to use cases where actions are sequentially generated while communicating with collaborators. Based on the above idea, we aimed to inductively express the relationship between motions and symbols contained in sensorimotor information~\cite{Harned1990pd} within a model.
Grounding each segmented subtask and language expression is necessary to express more primitive action combinations. The RNN expresses internal dynamics with the semantic and syntactic structure of the language system as a constraint by learning that the given language instructions and individual tasks are semantically grounded. This internal representation can be extended to cases where actions are sequentially generated while communicating with collaborators. Based on the above idea, we aim to inductively express the relationship between motions and symbols contained in sensorimotor information~\cite{Harned1990pd} within a model.



\textbf{Grounding of Language and Robot Motion:}
%In general, the correspondence between language, which is a discrete expression, and robot motions in a continuous space changes depending on the past context and the situation on the spot. For example, even with a simple instruction such as "Hit the blue bell," the action to accomplish this will vary depending on where the blue bell is placed, what posture the robot is in, etc. Robots need to acquire internal representations that can be used in a generalized manner even in unknown situations, so that they can achieve such ambiguous and one-time grounding each time.
In general, the correspondence between language, which is a discrete expression, and robot motion in a continuous space changes depending on the past context and situation on the spot. For example, even with simple instructions such as "Hit the bluebell," the action to accomplish this will vary depending on where the bluebell is placed, what position the robot is in, etc. Robots must acquire internal representations that can be generalized, even in unknown situations, to achieve such ambiguous and one-time grounding each time.


%Our DPL framework makes it possible to embed the understanding of language instructions and the motions into a set of attractor structures. \cite{Yamada2016fn,Yamada2017fn} acquire cyclic attractors that express the above episodic repetition as the internal dynamics by having RNN predict and learn a sequence that repeats multiple episodes consisting of waiting, verbal instructions, and robot motion generation. (Figure~\ref{fig_language_reduce}). When a verbal instruction and a camera image are input to the RNN at a fixed point attractor, the internal state of the RNN changes in a branching manner according to the local situation. After the language input is completed, robot motions are generated according to the meaning of the instructions in the environment. Then, once the motion is completed, the internal state converges to a fixed point again and waits for the next instruction. We demonstrated the generalization ability to generate appropriate actions for new situations without having to teach every combination of environmental conditions and verbal instructions.
Our DPL framework enables embedding an understanding of language instruction and motion into a set of attractor structures. \cite{Yamada2016fn,Yamada2017fn} acquire cyclic attractors that express the above episodic repetition as the internal dynamics by having RNN predict and learn a sequence that repeats multiple episodes of waiting, verbal instructions, and robot motion generation. (see Figure~\ref{fig_language_reduce}). When a verbal instruction and a camera image are input to the RNN at a fixed-point attractor, the internal state of the RNN changes in a branching manner according to the local situation. After the language input was completed, robot motions were generated according to the meaning of the instructions in the environment. Once the motion was complete, the internal state converged to a fixed point and waited for the next instruction. We demonstrated the generalization ability to generate appropriate actions for new situations without having to teach every combination of environmental conditions and verbal instructions.


%-------------------- Fig --------------------
\setlength\textfloatsep{5pt}
\begin{figure}[t]
\centering
\includegraphics[width=\columnwidth]{rsj2022_language_reduce.eps}
\caption{
Grounding of language instructions and motion generation in the designed internal dynamics of RNN~\cite{Yamada2016fn}.
}
\label{fig_language_reduce}
\hspace{10mm}
\end{figure}


%Furthermore, it is also possible to extend this model to a bidirectional conversion model that includes conversion from task instructions to actions and from actions to linguistic explanations. In \cite{Yamada2018ral}, we prepared two autoencoders, one for language and one for sensorimotor sequences, and performed learning by setting constraints such that the feature vectors of the corresponding sentences and motions were close to each other. After the training phase, the model achieved two-way translation from instructions to robot motions and robot motions to linguistic explanations via feature vectors. \cite{Toyoda2021ral,Toyoda2022ral} added that model a word distributed representation conversion module that utilizes the Word2vec model, allowing it to perform appropriate operations even for instructions that include new words similar to trained words.
Furthermore, extending this model to a bidirectional conversion model that includes conversion from task instructions to actions and from actions to linguistic explanations is possible. In \cite{Yamada2018ral}, we prepared two autoencoders, one for language and one for sensorimotor sequences. We performed learning by setting constraints such that the feature vectors of the corresponding sentences and motions were close to each other. After the training phase, the model achieved a two-way translation from instructions to robot motions and robot motions to linguistic explanations via feature vectors. \cite{Toyoda2021ral,Toyoda2022ral} added that model a word distributed representation conversion module that utilizes the Word2vec model, allowing it to perform appropriate operations even for instructions that include new words similar to trained words.


\textbf{Integration with Task Planning:}
%In the above studies, we introduced methods to give instructions in the form of language and integrate the robot's motions. Linguistic instructions are given word by word, but general-purpose planning description languages (e.g. Planning Domain Description Language: PDDL), which are used in conventional automatic planning methods, can be used as task instructions. \cite{Kase2020icra} proposed a learning method to connect descriptive language that defines the environment and images (Figure~\ref{fig_kase2022icra}). It designed motions using any automatic planning method based on discrete symbolic representations extracted from images. Although conventional PDDL makes it difficult to respond to changes in the environment because it is planned, our method makes it possible to dynamically determine the robot's actions.
In the above studies, we introduced methods to provide instructions in the language form and integrate the robot's motions. Linguistic instructions are given word by word; however, general-purpose planning description languages (e.g., planning domain description language: PDDL) used in conventional automatic planning methods can be used as task instructions. \cite{Kase2020icra} proposed a learning method to connect descriptive language that defines the environment and images (Figure~\ref{fig_kase2022icra}). They designed motions using an automatic planning method based on discrete symbolic representations extracted from images. Although conventional PDDL makes it difficult to respond to planned environmental changes, our method makes it possible to determine the robot's actions dynamically.


%In recent years, with the development of large language models (LLM~\cite{Zhao2023}), they have been actively applied to robot motion generation~\cite{Shridhar2021a,Liang2022a,Brohan2022a,Driess2023a}. Especially in task planning, it is possible to convert linguistic instruction to multiple subtask instructions in table format~\cite{Ahn2022,huang2022language} and code format~\cite{Vemprala2023,codeaspolicies2022}. Although our studies introduced in this section deal with limited language formats, we plan to expand them by incorporating LLM technology~\cite{Hori2024}.
In recent years, large language models (LLM~\cite{Zhao2023}) have been actively applied in robot motion generation~\cite{Shridhar2021a,Liang2022a,Brohan2022a,Driess2023a}. Especially in task planning, it is possible to convert linguistic instructions into multiple subtask instructions in table format~\cite{Ahn2022,huang2022language} and code format~\cite{Vemprala2023,codeaspolicies2022}. Although the studies introduced in this section deal with limited language formats, we plan to expand them by incorporating LLM technology~\cite{Hori2024}.


%-------------------- Fig --------------------
\begin{figure}[t]
\centering
\includegraphics[width=\columnwidth]{icra2020_overview.eps}
\caption{
Grounding with a generic planning language and extracted image features~\cite{Kase2020icra}.
}
\label{fig_kase2022icra}
\hspace{10mm}
\end{figure}    



\subsection{Application to Various Robot Tasks}
\subsubsection{Deformable Object Manipulation}

%We also verified the versatility of the DPL concept by applying it to difficult manipulation tasks. A typical example is the task of manipulating deformable objects. Manipulating with flexible objects such as cloth and strings that change shape upon contact makes it difficult to describe the entire situation in a program. \cite{Koma2016ral} generalized the motion of folding a towel by performing integrated learning based on image-motion information acquired through the robot's work experience. As a result, the robot could repeatedly perform tasks that correspond to object position shifts and different objects (such as books) from limited training data. Unlike conventional methods, our method does not require high-cost processing such as modeling or image processing, and it is possible to generate motion at a natural speed\footnote{https://youtu.be/YH1TrL1q6Po}. Our concept does not depend on tasks or robot hardware, so by implementing each function as a component, we easily reuse the implemented system in response to changes in requirements~\cite{Kanamura2021sensors}. In previous studies, we realized  the in-air knotting task~\cite{Suzuki2021iros} and the buttoning task~\cite{Fujii2022sii} by using the same model.
We verified the versatility of the DPL concept by applying it to difficult tasks. A typical example is the manipulation of deformable objects. Manipulating flexible objects such as cloths and strings that change shape upon contact makes it difficult to describe the entire situation in a program. \cite{Koma2016ral} generalized the motion of folding a towel by performing integrated learning based on image-motion information acquired through the robot's work experience. As a result, the robot could repeatedly perform tasks corresponding to object position shifts and different objects (such as books) using limited training data. Unlike conventional methods, our method does not require high-cost processing such as modeling or image processing, and it is possible to generate motion at a natural speed\footnote{https://youtu.be/YH1TrL1q6Po}. Our concept does not depend on tasks or robot hardware; therefore, by implementing each function as a component, we can easily reuse the implemented system in response to changes in requirements~\cite{Kanamura2021sensors}. We realize the in-air knotting task~\cite{Suzuki2021iros} and the buttoning task~\cite{Fujii2022sii} using the same model in previous studies.


%It has been confirmed that including detailed object information such as proximity sensors and tactile sensors as input to the model improves prediction performance. \cite{Ichiwara2022contact} introduced an attention mechanism that extracted information (such as the position and direction of the object) from the robot's visual and tactile information (Figure~\ref{fig_attention_bag}). In this task, it is necessary to generate an appropriate trajectory in real time according to the state of the bag and robot hand. The game starts with grabbing randomly placed zippers, so the gripping position and posture constantly change. The number of training sequences was 36, which is extremely small compared to general learning methods. In the experiment, it was confirmed that the perceptual predictions (predicted images, tactile states, and motions) of the model were corrected in real time by the attention mechanism.
It has been confirmed that including detailed object information, such as proximity and tactile sensors, as input to the model improves the prediction performance. \cite{Ichiwara2022contact} introduced an attention mechanism that extracted information (such as the position and direction of the object) from the robot's visual and tactile information (Figure~\ref{fig_attention_bag}). In this task, generating an appropriate trajectory in real-time is necessary according to the state of the bag and robot hand. The game starts by grabbing randomly placed zippers so that the gripping position and posture constantly change. The number of training sequences was 36, which is extremely small compared with general learning methods. The experiment confirmed that the model's perceptual predictions (predicted images, tactile states, and motions) were corrected in real time using an attention mechanism.


%The sense of touch also plays an important role in the emergence of diversity generated from exploration based on intrinsic motivation. Predictability and operability are criteria of the learning, however, they do not indicate the possibility of further exploration. We also proposed searching for new actions based on tactile perceptual information~\cite{Mori2020} and active inference based on the internal characteristics of the manipulated object~\cite{Saito2021,Saito2023a}. These studies addressed the next important issue to be solved in our DPL concept from the perspective of the emergence of various movements based on exploration.
Touch also plays an important role in the emergence of diversity generated from exploration based on intrinsic motivation. Predictability and operability are the criteria for learning; however, they do not indicate the possibility of further exploration. We also proposed searching for new actions based on tactile perceptual information~\cite{Mori2020} and active inference based on the internal characteristics of the manipulated object~\cite{Saito2021,Saito2023a}. These studies address the next important issue to be solved in our DPL concept from the perspective of the emergence of various movements based on exploration.


%-------------------- Fig --------------------
\setlength\textfloatsep{5pt}
\begin{figure}[t]
\centering
\includegraphics[width=\columnwidth]{attention_bag.eps}
\caption{
Generated bag opening task using attention mechanism with vision and tactile sensors~\cite{Ichiwara2022contact}.
}
\label{fig_attention_bag}
\hspace{10mm}
\end{figure}   



\subsubsection{Neuro-Robotics Experiments}

%We also studied to use a similar hierarchical predictive model as a robot simulation for neurodevelopmental disorders. We reproduced the functional disconnection between layers~\cite{idei2021paradoxical} and the imbalance of excitability and inhibition in neurons~\cite{idei2020homogeneous}. The results of these studies show that when a real robot causes an abnormality in sensory accuracy or even upper-layer prediction accuracy, this is observed as a behavioral abnormality. Although this paper does not go into detail, the above studies are significant in that they extend the DPL concept from a perspective not limited to task learning.
We also used a similar hierarchical predictive model for robot simulation of neurodevelopmental disorders. We reproduced the functional disconnection between layers~\cite{idei2021paradoxical} and the imbalance between excitability and inhibition in neurons~\cite{idei2020homogeneous}. These studies show that when a real robot causes an abnormality in sensory accuracy or even upper-layer prediction accuracy, it is observed as a behavioral abnormality. Although this paper does not go into detail, the studies mentioned above are significant because they extend the DPL concept from a perspective not limited to task learning.
